% unidades/06-pasiva-basica.tex
\chapter{\emoji{shield} \ruby{受け身}{うけみ} — Voz Pasiva Básica}
\unidadheader{06}{受け身}{受け身の基本}{Describir acciones recibidas o experimentadas}

\begin{tcolorbox}[vocabbox, title={\emoji{pushpin} Vocabulario}]
\begin{tabularx}{\linewidth}{llX}
\toprule
\textbf{Japonés} & \textbf{Español} & \textbf{Tipo} \\ \midrule
\vocabitem{褒める}{ほめる}{elogiar / halagar}{v. G2}
\vocabitem{叱る}{しかる}{regañar / reprender}{v. G1}
\vocabitem{噛む}{かむ}{morder}{v. G1}
\vocabitem{踏む}{ふむ}{pisar}{v. G1}
\vocabitem{盗む}{ぬすむ}{robar}{v. G1}
\vocabitem{誘う}{さそう}{invitar}{v. G1}
\vocabitem{選ぶ}{えらぶ}{elegir / seleccionar}{v. G1}
\bottomrule
\end{tabularx}
\end{tcolorbox}

\subsection*{\emoji{speech-balloon} Frases Clave}
\frase{\ruby{先生}{せんせい}に\ruby{褒}{ほ}められました。}{Fui elogiado/a por el profesor.}
\frase{\ruby{友達}{ともだち}に\ruby{誘}{さそ}われました。}{Fui invitado/a por mi amigo/a.}
\frase{\ruby{犬}{いぬ}に\ruby{噛}{か}まれた。}{Me mordió un perro.}
\frase{\ruby{電車}{でんしゃ}の\ruby{中}{なか}で\ruby{足}{あし}を\ruby{踏}{ふ}まれた。}{Me pisaron el pie en el tren.}

\subsection*{\emoji{gear} Gramática}
\patron{Formación de la voz pasiva}
  {G1: final う→われる \quad G2: る→られる \quad する→される \quad くる→こられる}
  {\ej{\ruby{書}{か}く → \ruby{書}{か}かれる。\ruby{読}{よ}む → \ruby{読}{よ}まれる。}{G1: la vocal final u→aれる}
   \ej{\ruby{食}{た}べる → \ruby{食}{た}べられる。\ruby{見}{み}る → \ruby{見}{み}られる。}{G2: る→られる}
   \ej{する → される。くる → こられる。}{G3: irregulares}}

\patron{Estructura: sujeto は/が [agente]に [V-pasiva]}
  {[主語]は [動作主]に [受け身動詞]}
  {\ej{\ruby{私}{わたし}は\ruby{先生}{せんせい}に\ruby{褒}{ほ}められました。}{Fui elogiado/a por el/la profesor/a.}
   \ej{\ruby{彼}{かれ}はクラスメートに\ruby{選}{えら}ばれた。}{Fue elegido por sus compañeros de clase.}
   \ej{その\ruby{提案}{ていあん}は\ruby{全員}{ぜんいん}に\ruby{支持}{しじ}された。}{Esa propuesta fue apoyada por todos.}}

\comparacion{Activa vs Pasiva}
  {\ruby{先生}{せんせい}が\ruby{私}{わたし}を\ruby{褒}{ほ}めた。\\\small La profesora me elogió.\\(enfoque: acción del agente)}
  {\ruby{私}{わたし}は\ruby{先生}{せんせい}に\ruby{褒}{ほ}められた。\\\small Fui elogiada por la profesora.\\(enfoque: experiencia del sujeto)}
  {La pasiva en japonés desplaza el foco hacia quien recibe o experimenta la acción. El agente se marca con に, no con が.}

\coloquial{〜られる (pasiva formal)}{〜れる (pasiva coloquial G1 / potencial)}{
  En habla casual, la distinción pasiva/potencial puede ambigüizarse en G2: 食べられる puede ser «puedo comer» o «fui comido». Contexto resuelve la ambigüedad.}

\culturabox{\ruby{受け身}{うけみ}\ruby{文化}{ぶんか} — La perspectiva de quien recibe}{
  En japonés, es muy común describir los eventos desde la perspectiva de quien los experimenta. En lugar de decir «el perro me mordió», el japonés usa la pasiva: «fui mordido por el perro». Esta perspectiva refleja una orientación hacia el impacto personal de los eventos.
}

\lectura{\ruby{大変}{たいへん}な\ruby{一日}{いちにち}\\[0.4em]
  \ruby{今日}{きょう}は\ruby{朝}{あさ}から\ruby{大変}{たいへん}だった。まず、\ruby{電車}{でんしゃ}の\ruby{中}{なか}で\ruby{足}{あし}を\ruby{踏}{ふ}まれた。会社では\ruby{上司}{じょうし}に\ruby{叱}{しか}られた。\ruby{夜}{よる}は\ruby{財布}{さいふ}を\ruby{盗}{ぬす}まれてしまった。\ruby{本当}{ほんとう}に\ruby{ひどい}{}\ruby{一日}{いちにち}だった。}
{Un día terrible\\[0.4em]
  Hoy fue terrible desde la mañana. Primero, me pisaron el pie en el tren. En el trabajo, me regañó mi jefe. Por la noche, me robaron la billetera. Realmente fue un día horrible.}

\begin{tcolorbox}[ejerciciobox, title={\emoji{pencil} Ejercicios}]
\begin{enumerate}
  \item Forma la voz pasiva: \ruby{褒}{ほ}める、\ruby{選}{えら}ぶ、する、\ruby{食}{た}べる
  \item Convierte a pasiva: \ruby{先生}{せんせい}が\ruby{学生}{がくせい}を\ruby{叱}{しか}った。
  \item ¿Activa o pasiva? ¿Cuándo preferirías usar la pasiva en este contexto?
\end{enumerate}
\textbf{Respuestas:}
1) \ruby{褒}{ほ}められる・\ruby{選}{えら}ばれる・される・\ruby{食}{た}べられる \quad
2) \ruby{学生}{がくせい}は\ruby{先生}{せんせい}に\ruby{叱}{しか}られた。\quad
3) (libre — reflexión sobre perspectiva)
\end{tcolorbox}

\begin{tcolorbox}[kanjibox, title={\emoji{mahjong} Kanjis}]
\begin{tabularx}{\linewidth}{cllXX}
\toprule
\textbf{Kanji} & \textbf{On} & \textbf{Kun} & \textbf{Significado} & \textbf{Vocab} \\ \midrule
\kanjirow{受}{じゅ}{\ruby{う}{ける}}{recibir}{\ruby{受}{う}ける / \ruby{受験}{じゅけん}}
\kanjirow{褒}{ほう}{\ruby{ほ}{める}}{elogiar}{\ruby{褒}{ほ}める / \ruby{褒美}{ほうび}}
\kanjirow{叱}{しつ}{\ruby{しか}{る}}{regañar}{\ruby{叱}{しか}る / \ruby{叱責}{しっせき}}
\kanjirow{選}{せん}{\ruby{えら}{ぶ}}{elegir}{\ruby{選}{えら}ぶ / \ruby{選択}{せんたく}}
\bottomrule
\end{tabularx}
\end{tcolorbox}
